\documentclass[aspectratio=169]{beamer}
\usepackage{iftex}
\ifPDFTeX\usepackage[T1]{fontenc}\usepackage{lmodern}\fi
\usepackage[czech]{babel}
\usepackage{listings}
\newcommand{\CVUTTemplatePath}{../../../utils/presentation-template}
\usepackage{\CVUTTemplatePath/beamerthemeCVUT}
\setbeamertemplate{footline}[frame number]
\beamertemplatenavigationsymbolsempty
\newcommand{\LANG}{CZ}\newcommand{\FRONTPAGE}{dark}
\title{Základní principy testování}
\subtitle{Základy programování\\Matouš Cejnek\\U12110, FS, ČVUT v Praze}
\author{}\institute{}\date{}
\renewcommand{\headlineA}{}\renewcommand{\headlineB}{}\renewcommand{\headlineC}{}
\definecolor{CodeBlue}{RGB}{38,83,125}\definecolor{SoftGray}{RGB}{242,244,247}
\definecolor{CodeGreen}{RGB}{38,125,94}\definecolor{WarningRed}{RGB}{170,45,45}
\setbeamerfont{normal text}{size=\large}\setbeamerfont{frametitle}{size=\LARGE}
\setbeamercolor{block title}{fg=white,bg=CodeBlue}
\setbeamercolor{block body}{fg=black,bg=SoftGray}
\newcommand{\term}[1]{\textcolor{CodeBlue}{\textbf{#1}}}
\lstdefinestyle{python}{language=Python,basicstyle=\ttfamily\normalsize,
 keywordstyle=\color{CodeBlue}\bfseries,stringstyle=\color{CodeGreen},
 commentstyle=\color{gray},showstringspaces=false,columns=fullflexible,keepspaces=true}
\begin{document}
\begin{frame}\clearpage\thispagestyle{empty}\titlepage\end{frame}

\begin{frame}{Co si dnes odnesete?}
\begin{itemize}
\item Vysvětlit, co test ověřuje a co naopak nezaručuje.
\item Psát malé testy pomocí \texttt{assert} a frameworku pytest.
\item Navrhovat běžné, hraniční i chybové případy.
\item Testovat výjimky, desetinná čísla a stav objektů.
\item Používat testy jako spustitelnou specifikaci.
\end{itemize}
\end{frame}

\begin{frame}{Proč nestačí program jednou spustit?}
\begin{itemize}
\item Jeden běh pokryje jen jednu cestu programem.
\item Ruční kontrola se špatně opakuje po každé změně.
\item Chyba se často projeví až na hranici nebo neobvyklém vstupu.
\item Výstup, který vypadá rozumně, nemusí být správný.
\end{itemize}
\begin{block}{Cíl testování}Rychle získat důvěryhodnou informaci o konkrétním očekávaném chování.\end{block}
\end{frame}

\begin{frame}{Test je příklad i kontrola}
\begin{itemize}
\item Připraví vstup a počáteční stav.
\item Spustí jednu část programu.
\item Porovná skutečný výsledek s očekáváním.
\item Selže čitelným způsobem, pokud se chování liší.
\end{itemize}
\end{frame}

\begin{frame}[fragile]{Nejmenší kontrola: \texttt{assert}}
\begin{lstlisting}[style=python]
def area(width, height):
    return width * height

assert area(3, 4) == 12
assert area(0, 4) == 0
\end{lstlisting}
\begin{block}{Význam}Pokud je výraz nepravdivý, Python vyvolá \texttt{AssertionError}.\end{block}
\end{frame}

\begin{frame}{Co patří do testu?}
\begin{description}
\item[Arrange] Připrav data, objekty a závislosti.
\item[Act] Proveď právě testovanou operaci.
\item[Assert] Ověř pozorovatelný výsledek.
\end{description}
\begin{block}{Arrange--Act--Assert}Struktura pomáhá rychle poznat, proč test selhal.\end{block}
\end{frame}

\begin{frame}[fragile]{Testovací funkce v pytestu}
\begin{lstlisting}[style=python]
from geometry import area

def test_area_of_rectangle():
    width, height = 3, 4       # Arrange
    result = area(width, height)  # Act
    assert result == 12        # Assert
\end{lstlisting}
\end{frame}

\begin{frame}[fragile]{Jak pytest testy najde?}
\begin{itemize}
\item Soubory obvykle pojmenujeme \texttt{test\_*.py}.
\item Testovací funkce začínají \texttt{test\_}.
\item Příkaz \texttt{pytest} provede sběr a spuštění testů.
\item Výstup ukáže místo selhání i porovnávané hodnoty.
\end{itemize}
\begin{lstlisting}[basicstyle=\ttfamily\large]
python -m pytest
\end{lstlisting}
\end{frame}

\begin{frame}{Testovací sada není důkaz bezchybnosti}
\begin{itemize}
\item Úspěch znamená jen: prošly všechny napsané testy.
\item Nenapsaný případ nebyl ověřen.
\item Chybný test může potvrzovat chybné očekávání.
\item Důležitější než počet testů je volba případů.
\end{itemize}
\end{frame}

\begin{frame}{Tři skupiny vstupů}
\begin{description}
\item[Běžné] Reprezentativní hodnoty z normálního provozu.
\item[Hraniční] Nula, minimum, maximum, prázdná kolekce, přechod pásem.
\item[Chybové] Neplatný typ, formát nebo hodnota.
\end{description}
\end{frame}

\begin{frame}{Hranice rozděluje chování}
\begin{itemize}
\item Pro podmínku \texttt{x < 10} testujeme 9, 10 a 11.
\item Pro interval testujeme oba kraje i hodnoty těsně vedle.
\item U kolekce testujeme prázdnou, jednoprvkovou a vícepoložkovou.
\item U textu testujeme prázdný řetězec a zvláštní znaky.
\end{itemize}
\end{frame}

\begin{frame}[fragile]{Výjimka může být správný výsledek}
\begin{lstlisting}[style=python]
import pytest

def test_division_by_zero_is_rejected():
    with pytest.raises(ZeroDivisionError):
        divide(10, 0)
\end{lstlisting}
\begin{block}{Pozor}Test má ověřit konkrétní očekávanou výjimku, ne libovolnou chybu.\end{block}
\end{frame}

\begin{frame}[fragile]{Desetinná čísla porovnáváme s tolerancí}
\begin{lstlisting}[style=python]
import pytest

def test_average():
    result = average([0.1, 0.2])
    assert result == pytest.approx(0.15)
\end{lstlisting}
\begin{itemize}
\item Binární reprezentace nemusí být matematicky přesná.
\item Tolerance musí odpovídat významu výpočtu.
\end{itemize}
\end{frame}

\begin{frame}{Jeden test, jedna hlavní myšlenka}
\begin{itemize}
\item Název testu popisuje očekávané chování.
\item Při selhání má být zřejmé, co se porušilo.
\item Nesouvisející scénáře patří do samostatných testů.
\item Několik kontrol jednoho výsledku může být v jednom testu.
\end{itemize}
\end{frame}

\begin{frame}{Testujeme rozhraní, ne vnitřnosti}
\begin{itemize}
\item Ověřujeme návratovou hodnotu, výjimku nebo veřejný stav.
\item Soukromé pomocné kroky se mohou při refaktoringu změnit.
\item Test nemá bránit změně implementace se stejným chováním.
\end{itemize}
\begin{block}{Dobrá otázka}Co může uživatel funkce nebo objektu pozorovat?\end{block}
\end{frame}

\begin{frame}[fragile]{Test objektu sleduje změnu stavu}
\begin{lstlisting}[style=python]
def test_deposit_increases_balance():
    account = Account(balance=100)
    account.deposit(50)
    assert account.balance == 150
\end{lstlisting}
\begin{itemize}\item Každý test vytváří vlastní čerstvý objekt.\end{itemize}
\end{frame}

\begin{frame}{Izolace testů}
\begin{itemize}
\item Pořadí spuštění nemá ovlivnit výsledek.
\item Testy nemají sdílet měnitelný globální stav.
\item Dočasná data se připraví a po testu uklidí.
\item Selhání jednoho testu nemá rozbít další testy.
\end{itemize}
\end{frame}

\begin{frame}{Fixture připravuje opakovaný kontext}
\begin{itemize}
\item Fixture poskytuje testovací data nebo objekt.
\item Pytest ji předá testu podle názvu parametru.
\item Snižuje duplicitu, ale nesmí skrývat podstatný scénář.
\item Vhodná je například dočasná složka \texttt{tmp\_path}.
\end{itemize}
\end{frame}

\begin{frame}{Parametrizace: stejné pravidlo, více příkladů}
\begin{itemize}
\item Jeden testovací postup dostane několik dvojic vstup--výstup.
\item Každý případ je ve výsledku veden samostatně.
\item Tabulka případů zpřehlední hranice a ekvivalenční třídy.
\item Nepoužívejme parametrizaci pro zásadně odlišné scénáře.
\end{itemize}
\end{frame}

\begin{frame}{Regresní test zachytí nalezenou chybu}
\begin{enumerate}
\item Chybu nejprve reprodukujeme testem.
\item Ověříme, že test na starém kódu selže.
\item Opravíme implementaci.
\item Test ponecháme jako ochranu před návratem chyby.
\end{enumerate}
\end{frame}

\begin{frame}{Testy a změna programu}
\begin{itemize}
\item Refaktoring mění strukturu, ale zachovává chování.
\item Procházející testy dávají rychlou zpětnou vazbu.
\item Změna požadavku obvykle vyžaduje změnu testu i kódu.
\item Testovací sada je živá součást programu.
\end{itemize}
\end{frame}

\begin{frame}{Test jako specifikace pro AI}
\begin{itemize}
\item Příklady zpřesňují slovní zadání.
\item Hraniční případy omezují nežádoucí interpretace.
\item Vygenerovaný kód se hodnotí podle chování, ne dojmu.
\item Testy vytvořené stejným modelem je nutné kriticky zkontrolovat.
\end{itemize}
\end{frame}

\begin{frame}{Časté chyby v testech}
\begin{itemize}
\item Test jen zopakuje implementovaný výpočet.
\item Chybí hranice a chybové stavy.
\item Test závisí na čase, síti nebo pořadí spuštění.
\item Název neříká, jaké chování se očekává.
\item Příliš mnoho scénářů způsobí nečitelný neúspěch.
\end{itemize}
\end{frame}

\begin{frame}{Pracovní postup}
\begin{enumerate}
\item Popište očekávané chování konkrétním příkladem.
\item Zvolte běžné, hraniční a chybové případy.
\item Napište malý čitelný test a spusťte ho.
\item Opravte příčinu, nikoli pouze příznak selhání.
\item Spusťte celou sadu po každé změně.
\end{enumerate}
\end{frame}

\begin{frame}{Odkud vezmeme očekávaný výsledek?}
\begin{itemize}
\item Ze zadání, dokumentace nebo doménového pravidla.
\item Z ručně spočítaného jednoduchého příkladu.
\item Z nezávislé referenční implementace.
\item Z vlastnosti, která musí platit pro mnoho vstupů.
\end{itemize}
\begin{block}{Testovací orákulum}Zdroj, podle kterého rozhodneme, zda je výsledek správný.\end{block}
\end{frame}

\begin{frame}{Ekvivalenční třídy vstupů}
\begin{itemize}
\item Vstupy rozdělíme podle očekávaného chování.
\item Z každé skupiny vybereme reprezentativní příklad.
\item Hranice mezi skupinami ověříme zvlášť.
\item Nezkoušíme všechny hodnoty, ale systematicky pokrýváme pravidla.
\end{itemize}
\end{frame}

\begin{frame}{Úrovně testování}
\begin{description}
\item[Jednotkový test] Ověřuje malou část programu izolovaně.
\item[Integrační test] Ověřuje spolupráci několika částí.
\item[Systémový test] Ověřuje celý pracovní postup jako celek.
\end{description}
\begin{itemize}\item Různé úrovně hledají různé druhy problémů.\end{itemize}
\end{frame}

\begin{frame}{Rychlost a rozsah testů}
\begin{itemize}
\item Malých rychlých testů může běžet mnoho po každé změně.
\item Integrační testy jsou pomalejší, ale ověřují skutečné propojení.
\item Několik systémových testů chrání nejdůležitější scénáře.
\item Pomalý test se spouští méně často a dává pozdější zpětnou vazbu.
\end{itemize}
\end{frame}

\begin{frame}{Deterministický test}
\begin{itemize}
\item Za stejných podmínek vždy dopadne stejně.
\item Náhodnost řídíme pevným seedem nebo ji předáme jako závislost.
\item Aktuální čas nepoužíváme skrytě uvnitř testované logiky.
\item Síťové služby a cizí data mohou způsobit nestabilní testy.
\end{itemize}
\end{frame}

\begin{frame}{Pokrytí kódu je mapa, ne cíl}
\begin{itemize}
\item Pokrytí ukazuje, které řádky či větve testy navštívily.
\item Navštívený řádek ještě nemusí být dobře ověřen.
\item Nízké pokrytí odhalí zjevně netestované oblasti.
\item Vysoké procento nenahrazuje kvalitní očekávání a případy.
\end{itemize}
\end{frame}

\begin{frame}{Když test selže}
\begin{enumerate}
\item Přečteme první relevantní chybovou zprávu.
\item Ověříme, zda je správný test i jeho očekávání.
\item Zmenšíme problém na nejjednodušší reprodukci.
\item Opravíme příčinu a spustíme nejprve dotčený test.
\item Nakonec spustíme celou sadu kvůli regresím.
\end{enumerate}
\end{frame}

\begin{frame}{Shrnutí}
\begin{itemize}
\item Test je opakovatelná kontrola očekávaného chování.
\item Arrange--Act--Assert odděluje přípravu, akci a kontrolu.
\item Kvalitu sady určují především dobře zvolené případy.
\item Pytest podporuje výjimky, aproximace, fixtures i parametrizaci.
\item Úspěšné testy zvyšují důvěru, ale nejsou důkazem bezchybnosti.
\end{itemize}
\end{frame}
\end{document}
